% =====================================================================
%  CHƯƠNG 2: CƠ SỞ LÝ THUYẾT
% =====================================================================
\chapter{Cơ sở lý thuyết}
\label{chap:background}

\section{Hệ thống gợi ý theo hướng truy hồi}

Gọi tập người dùng là $U$ và tập sản phẩm là $I$. Thay vì dự đoán điểm số cho từng
cặp như lọc cộng tác cổ điển, đồ án mô hình hóa bài toán theo hướng \textbf{truy hồi
dựa trên biểu diễn} (embedding-based retrieval): mỗi người dùng $u$ và mỗi sản phẩm
$i$ được ánh xạ thành các vector $\mathbf{u}, \mathbf{v} \in \mathbb{R}^d$ trong cùng
một không gian. Mức độ phù hợp được đo bằng độ tương đồng cosine:

\begin{equation}
    \text{score}(u, i) = \cos(\mathbf{u}, \mathbf{v})
    = \frac{\mathbf{u} \cdot \mathbf{v}}{\|\mathbf{u}\|\,\|\mathbf{v}\|}.
\end{equation}

Khi cả hai vector đã được chuẩn hóa $L_2$ (độ dài bằng $1$), độ tương đồng cosine
trở thành tích vô hướng đơn giản. Gợi ý Top-$N$ chính là việc tìm $N$ sản phẩm có
score cao nhất với người dùng -- một bài toán \textit{tìm kiếm láng giềng gần nhất}
(nearest neighbor search) trong không gian vector, có thể tăng tốc bằng các thư viện
chuyên dụng. Hướng tiếp cận này là nền tảng của các hệ thống gợi ý công nghiệp quy
mô lớn nhờ khả năng tách biệt việc tính toán biểu diễn người dùng và sản phẩm.

\section{Các hướng nghiên cứu liên quan}

Lĩnh vực hệ thống gợi ý đã trải qua nhiều giai đoạn phát triển:

\begin{itemize}[leftmargin=2em, itemsep=2pt]
    \item \textbf{Lọc cộng tác cổ điển:} từ các phương pháp dựa trên láng giềng
    (user/item-based KNN) tới \textit{phân rã ma trận} (Matrix Factorization, ALS),
    học các nhân tố ẩn từ ma trận tương tác. Hạn chế chung là chỉ dùng tín hiệu định
    danh và gặp khó với cold-start.
    \item \textbf{Học sâu cho gợi ý:} các mô hình như Neural Collaborative Filtering
    thay tích vô hướng bằng mạng nơ-ron để học tương tác phi tuyến, song vẫn dựa trên
    embedding theo định danh.
    \item \textbf{Truy hồi hai tháp:} được phổ biến bởi các hệ thống công nghiệp quy
    mô lớn (YouTube, Google), kiến trúc hai tháp cho phép truy hồi ứng viên cực nhanh
    từ kho hàng triệu mục, thường huấn luyện bằng học tương phản với mẫu âm trong
    batch.
    \item \textbf{Gợi ý dựa trên ngôn ngữ:} xu hướng gần đây kết hợp mô hình ngôn
    ngữ để khai thác văn bản mô tả sản phẩm, giúp biểu diễn giàu ngữ nghĩa và giảm
    phụ thuộc vào lịch sử tương tác -- đặc biệt hữu ích cho cold-start.
\end{itemize}

Đồ án kế thừa và kết hợp hai hướng cuối: dùng \textit{mô hình ngôn ngữ} làm tháp sản
phẩm trong kiến trúc \textit{hai tháp}, huấn luyện bằng \textit{học tương phản} --
một thiết kế hiện đại, phù hợp với bối cảnh dữ liệu lớn và đặc thù giàu văn bản của
sản phẩm phần mềm.

\section{Nền tảng dữ liệu lớn}

\subsection{Apache Spark}

\textbf{Apache Spark} là nền tảng tính toán phân tán cho phép xử lý dữ liệu quy mô
lớn theo cơ chế song song trong bộ nhớ. Trong đồ án, Spark đảm nhận toàn bộ khâu
tiền xử lý nặng: nạp dữ liệu JSON/JSONL thô, khử trùng lặp, ghép nối (join) các tập
dữ liệu và biến đổi đặc trưng. API \texttt{DataFrame} cùng bộ tối ưu truy vấn
\textit{Catalyst} giúp các phép biến đổi được thực thi hiệu quả; kỹ thuật
\textbf{broadcast join} được dùng để ghép tập tương tác nhỏ với bảng dữ liệu lớn mà
không gây xáo trộn (shuffle) tốn kém.

\subsection{Định dạng cột Apache Parquet}

Sau tiền xử lý, dữ liệu được lưu dưới định dạng \textbf{Apache Parquet} -- một định
dạng lưu trữ \textit{theo cột} (columnar). So với định dạng dòng (CSV, JSON),
Parquet mang lại các ưu điểm then chốt cho dữ liệu lớn:

\begin{itemize}[leftmargin=2em, itemsep=2pt]
    \item \textbf{Nén hiệu quả:} dữ liệu cùng cột có kiểu giống nhau nên nén tốt,
    giảm đáng kể dung lượng lưu trữ.
    \item \textbf{Đọc chọn lọc:} chỉ cần đọc các cột liên quan, giảm I/O.
    \item \textbf{Tốc độ đọc/ghi cao:} phù hợp với các tác vụ phân tích và huấn
    luyện lặp lại.
\end{itemize}

Để tăng tốc độ nạp dữ liệu khi huấn luyện, đồ án áp dụng kỹ thuật \textbf{nạp dữ
liệu đa luồng bất đồng bộ} (asynchronous multi-worker loading) thông qua
\texttt{DataLoader}, cho phép phần cứng CPU chuẩn bị batch tiếp theo trong khi GPU
đang tính toán batch hiện tại.

\section{Mô hình ngôn ngữ và biểu diễn văn bản}

\subsection{Kiến trúc Transformer}

\textbf{Mô hình ngôn ngữ} hiện đại dựa trên kiến trúc \textbf{Transformer} với cơ
chế \textit{tự chú ý} (self-attention), cho phép mô hình nắm bắt mối quan hệ ngữ
nghĩa giữa các từ trong văn bản bất kể khoảng cách. Cơ chế chú ý tính trọng số liên
kết giữa các token dựa trên ba ma trận truy vấn (Query), khóa (Key) và giá trị
(Value):

\begin{equation}
    \text{Attention}(Q, K, V) = \text{softmax}\!\left(\frac{Q K^\top}{\sqrt{d_k}}\right) V,
\end{equation}

trong đó $d_k$ là số chiều của vector khóa, đóng vai trò hệ số tỷ lệ giúp ổn định
gradient. Mỗi tầng Transformer (Hình~\ref{fig:encoder}) gồm một khối \textit{chú ý
đa đầu} (multi-head attention) và một mạng \textit{truyền thẳng} (feed-forward),
mỗi khối kèm kết nối tắt (residual) và chuẩn hóa lớp (layer norm). MiniLM xếp chồng
$6$ tầng như vậy. Nhờ tính toán song song trên toàn bộ chuỗi, Transformer học được
biểu diễn ngữ cảnh phong phú.

\begin{figure}[H]
\centering
\begin{tikzpicture}[
    font=\scriptsize, node distance=0.45cm,
    blk/.style={rectangle, rounded corners=2pt, draw=blue!55, fill=blue!8,
        line width=0.7pt, text width=3.4cm, minimum height=0.6cm, align=center},
    an/.style={rectangle, rounded corners=2pt, draw=teal!55, fill=teal!10,
        line width=0.7pt, text width=3.4cm, minimum height=0.5cm, align=center},
    a/.style={-{Stealth[length=1.8mm]}, line width=0.6pt, blue!55}
]
\node[blk, fill=gray!10, draw=gray!55] (in) {Embedding token + vị trí};
\node[blk, above=of in] (mha) {Multi-Head\\Self-Attention};
\node[an, above=of mha] (an1) {Add \& LayerNorm};
\node[blk, above=of an1] (ffn) {Feed-Forward (FFN)};
\node[an, above=of ffn] (an2) {Add \& LayerNorm};
\node[blk, fill=gray!10, draw=gray!55, above=of an2] (out) {Biểu diễn ngữ cảnh};
\draw[a] (in)--(mha); \draw[a] (mha)--(an1); \draw[a] (an1)--(ffn);
\draw[a] (ffn)--(an2); \draw[a] (an2)--(out);
% residual arrows
\draw[a, gray!55, dashed] (in.east) -- ++(1.3,0) |- (an1.east);
\draw[a, gray!55, dashed] (an1.west) -- ++(-1.3,0) |- (an2.west);
\node[draw=blue!45, rounded corners=3pt, dashed, fit=(mha)(an1)(ffn)(an2),
      inner sep=5pt, label={[font=\scriptsize\bfseries, blue!55!black]right:$\times\,6$ tầng}] {};
\end{tikzpicture}
\caption{Cấu trúc một tầng mã hóa Transformer (MiniLM xếp chồng $6$ tầng)}
\label{fig:encoder}
\end{figure} Đồ án sử dụng mô hình \texttt{all-MiniLM-L6-v2} thuộc họ
\textit{Sentence-Transformers} -- một mô hình gọn nhẹ ($6$ tầng, đầu ra $384$ chiều,
khoảng $22{,}7$ triệu tham số) được chưng cất (distillation) và tinh chỉnh chuyên để
sinh \textbf{vector câu} (sentence embedding) chất lượng cao, cân bằng tốt giữa độ
chính xác và tốc độ -- phù hợp với yêu cầu mã hóa hàng trăm nghìn văn bản.

\subsection{Từ văn bản tới vector ngữ nghĩa}

Văn bản mô tả sản phẩm (tiêu đề, tính năng, mô tả, danh mục\ldots) được đưa qua mô
hình để thu các vector ngữ cảnh cho từng token. Để có một vector \textit{cố định}
đại diện cho toàn bộ văn bản, đồ án áp dụng phép \textbf{trung bình có mặt nạ}
(mean pooling) trên các token:

\begin{equation}
    \mathbf{v} = \frac{\sum_{t} m_t \, \mathbf{h}_t}{\sum_{t} m_t},
\end{equation}

với $\mathbf{h}_t$ là vector của token thứ $t$ và $m_t \in \{0, 1\}$ là mặt nạ chú ý
(loại bỏ token đệm). Vector kết quả được chuẩn hóa $L_2$. Nhờ vậy, hai sản phẩm có
mô tả tương đồng về ngữ nghĩa sẽ có vector gần nhau, giải quyết được vấn đề khởi động
lạnh cho sản phẩm mới.

\subsection{Tiền huấn luyện và Sentence-BERT}

Sức mạnh của bộ mã hóa đến từ quá trình \textbf{tiền huấn luyện} (pre-training) trên
kho văn bản khổng lồ, giúp mô hình tích lũy tri thức ngôn ngữ tổng quát trước khi áp
dụng cho bài toán cụ thể. Tuy nhiên, các vector token từ BERT gốc không tối ưu cho
việc đo độ tương đồng câu. \textbf{Sentence-BERT} khắc phục bằng cách tinh chỉnh mô
hình theo cặp/bộ ba câu với mục tiêu sao cho các câu cùng nghĩa có vector gần nhau --
chính là tinh thần của học tương phản. Mô hình \texttt{all-MiniLM-L6-v2} được huấn
luyện trên hơn một tỷ cặp câu, nên vector $384$ chiều mà nó sinh ra mang chất lượng
ngữ nghĩa cao, rất phù hợp để làm \textit{biểu diễn sản phẩm cố định} trong tháp Item.
Đây là lý do đồ án có thể \textbf{đóng băng} tháp này mà vẫn thu được không gian biểu
diễn tốt, tiết kiệm đáng kể tài nguyên huấn luyện.

\section{Kiến trúc hai tháp (Two-Tower)}
\label{sec:twotower}

\textbf{Two-Tower} (hay \textit{dual encoder}) là kiến trúc gồm hai nhánh mạng độc
lập (Hình~\ref{fig:twotower}): một \textbf{tháp sản phẩm} mã hóa sản phẩm và một
\textbf{tháp người dùng} mã hóa người dùng, cùng cho ra vector trong một không gian
chung. Điểm mạnh của kiến trúc này là hai tháp có thể tính toán \textit{tách rời}:
toàn bộ vector sản phẩm có thể được mã hóa trước (offline) và lập chỉ mục, còn vector
người dùng chỉ cần tính khi cần gợi ý -- nhờ vậy hệ thống truy hồi cực nhanh trên kho
sản phẩm lớn.

\begin{figure}[H]
\centering
\begin{tikzpicture}[
    node distance=0.5cm and 1.0cm,
    box/.style={rectangle, rounded corners=3pt, draw=blue!55, fill=blue!9,
        line width=0.8pt, text width=3.0cm, minimum height=0.8cm, font=\footnotesize},
    frozen/.style={rectangle, rounded corners=3pt, draw=gray!60, fill=gray!12,
        line width=0.8pt, text width=3.0cm, minimum height=0.8cm, font=\footnotesize},
    op/.style={circle, draw=teal!65, fill=teal!12, line width=0.9pt,
        minimum size=0.95cm, font=\small},
    arr/.style={-{Stealth[length=2mm]}, line width=0.7pt, blue!60}
]
% User tower (left)
\node[box] (uin) {Vector lịch sử người dùng (384-d)};
\node[box, above=of uin] (ump) {Tháp User: MLP\\$384{\to}512{\to}384$\\BatchNorm + Dropout};
\node[box, above=of ump] (unorm) {Chuẩn hóa $L_2$};
\draw[arr] (uin) -- (ump);
\draw[arr] (ump) -- (unorm);
% Item tower (right)
\node[frozen, right=2.4cm of uin] (iin) {Văn bản sản phẩm\\(semantic\_text)};
\node[frozen, above=of iin] (ienc) {Tháp Item: MiniLM\\(\textbf{đóng băng})\\+ mean pooling};
\node[frozen, above=of ienc] (inorm) {Chuẩn hóa $L_2$};
\draw[arr, gray!60] (iin) -- (ienc);
\draw[arr, gray!60] (ienc) -- (inorm);
% similarity
\node[op] (sim) at ($(unorm)!0.5!(inorm) + (0,1.2)$) {$\cos$};
\draw[arr] (unorm) -- (sim);
\draw[arr, gray!60] (inorm) -- (sim);
\node[above=0.35cm of sim, font=\footnotesize\bfseries, text=teal!50!black]
    {Độ tương đồng → Loss InfoNCE};
\end{tikzpicture}
\caption{Kiến trúc hai tháp (Two-Tower) với tháp sản phẩm đóng băng}
\label{fig:twotower}
\end{figure}

Trong thiết kế của đồ án, \textbf{tháp sản phẩm được đóng băng hoàn toàn} (không cập
nhật trọng số) nhằm tiết kiệm bộ nhớ và tận dụng tri thức ngôn ngữ sẵn có; chỉ
\textbf{tháp người dùng} (một mạng nơ-ron đa tầng) được huấn luyện để học cách ánh
xạ lịch sử hành vi sang không gian sản phẩm.

\section{Học tương phản và hàm mất mát InfoNCE}
\label{sec:infonce}

\textbf{Học tương phản} (Contrastive Learning) huấn luyện mô hình bằng cách so sánh:
với mỗi người dùng, sản phẩm thực sự tương tác (mẫu dương) cần được kéo lại gần,
trong khi các sản phẩm khác (mẫu âm) bị đẩy ra xa. Đồ án sử dụng hàm mất mát
\textbf{InfoNCE} với chiến lược \textit{mẫu âm trong batch} (in-batch negatives):
trong một batch gồm $B$ cặp (người dùng, sản phẩm), sản phẩm của các cặp còn lại
đóng vai trò mẫu âm cho người dùng đang xét. Hàm mất mát cho người dùng $i$ là:

\begin{equation}
    \mathcal{L}_i = -\log
    \frac{\exp\!\big(\cos(\mathbf{u}_i, \mathbf{v}_i)/\tau\big)}
         {\sum_{j=1}^{B} \exp\!\big(\cos(\mathbf{u}_i, \mathbf{v}_j)/\tau\big)},
\end{equation}

trong đó $\mathbf{v}_i$ là sản phẩm dương của người dùng $i$, $\mathbf{v}_j$ là các
sản phẩm trong batch, và $\tau$ là \textit{tham số nhiệt độ} (temperature) điều
chỉnh độ "sắc" của phân phối. Về bản chất, đây là một bài toán phân loại $B$ lớp mà
nhãn đúng nằm trên đường chéo của ma trận tương đồng -- cài đặt hiệu quả bằng
\texttt{CrossEntropyLoss}. Chiến lược in-batch negatives giúp tận dụng tối đa mỗi
batch mà không cần lấy mẫu âm tường minh, rất phù hợp với huấn luyện quy mô lớn.
Hình~\ref{fig:infonce} minh họa trực quan: ma trận tương đồng $B \times B$ giữa mọi
người dùng và mọi sản phẩm trong batch, trong đó nhiệm vụ huấn luyện là làm cho các
ô \textit{đường chéo} (cặp đúng) có giá trị cao nhất trên mỗi hàng.

\begin{figure}[H]
\centering
\includegraphics[width=0.58\textwidth]{infonce_matrix.png}
\caption{Ma trận tương đồng in-batch của InfoNCE; đường chéo là cặp dương cần tối đa hóa}
\label{fig:infonce}
\end{figure}

\subsection{Vai trò của tham số nhiệt độ}

Tham số nhiệt độ $\tau$ kiểm soát mức độ "gay gắt" của hàm mất mát. Khi $\tau$ nhỏ,
mô hình \textit{trừng phạt mạnh} những mẫu âm khó (nằm gần mẫu dương), thúc đẩy phân
tách rạch ròi trong không gian biểu diễn; khi $\tau$ lớn, gradient trải đều và mô
hình "khoan dung" hơn. Về mặt lý thuyết, học tương phản tối ưu đồng thời hai tính
chất: \textbf{alignment} (các cặp dương phải gần nhau) và \textbf{uniformity} (các
biểu diễn phải trải đều trên mặt cầu đơn vị để giữ nhiều thông tin). Giá trị
$\tau = 0{,}1$ được đồ án lựa chọn là một mức cân bằng phổ biến và hiệu quả cho bài
toán gợi ý.

\section{Tìm kiếm láng giềng gần nhất}
\label{sec:ann}

Sau khi huấn luyện, việc gợi ý Top-$N$ quy về tìm các sản phẩm có vector gần vector
người dùng nhất. Với kho chỉ vài chục nghìn sản phẩm, ta có thể tính tương đồng
\textit{vét cạn} (brute-force) như trong giao thức đánh giá của đồ án. Tuy nhiên, khi
kho lên tới hàng triệu sản phẩm, phép vét cạn trở nên quá tốn kém cho yêu cầu thời
gian thực. Khi đó, các kỹ thuật \textbf{tìm kiếm láng giềng gần nhất xấp xỉ}
(Approximate Nearest Neighbor -- ANN) như chỉ mục đảo (IVF), lượng tử hóa tích
(product quantization) hay đồ thị HNSW -- được hiện thực trong thư viện
\textbf{FAISS} -- cho phép truy hồi trong vài phần nghìn giây với độ chính xác gần
như không suy giảm. Kiến trúc hai tháp đặc biệt phù hợp với ANN vì toàn bộ vector
sản phẩm có thể được lập chỉ mục trước. Đây là nền tảng cho hướng phát triển triển
khai thời gian thực ở Chương~\ref{chap:conclusion}.

\section{Các độ đo đánh giá}
\label{sec:metrics}

Vì bài toán là truy hồi Top-$N$, đồ án dùng hai độ đo xếp hạng phổ biến, tính trên
\textbf{toàn bộ kho sản phẩm}. Với mỗi người dùng, gọi $\text{Rel}_u$ là tập sản
phẩm thực sự liên quan và $\text{Rec}_u^K$ là Top-$K$ sản phẩm được gợi ý:

\begin{equation}
    \text{Recall@K} = \frac{|\text{Rec}_u^K \cap \text{Rel}_u|}{|\text{Rel}_u|}.
\end{equation}

Trong giao thức leave-one-out (mỗi người dùng có đúng một sản phẩm mục tiêu),
Recall@K rút gọn thành tỷ lệ người dùng có sản phẩm mục tiêu nằm trong Top-$K$.
\textbf{NDCG@K} bổ sung yếu tố thứ hạng, thưởng nhiều hơn khi sản phẩm đúng xuất hiện
ở vị trí cao:

\begin{equation}
    \text{NDCG@K} = \frac{1}{\log_2(p + 1)},
\end{equation}

với $p$ là vị trí của sản phẩm mục tiêu trong danh sách gợi ý (nếu nằm trong Top-$K$,
ngược lại bằng $0$). Việc đánh giá trên toàn kho $\sim 17.000$ sản phẩm khiến các
con số tuyệt đối nhỏ, nhưng phản ánh \textit{trung thực} độ khó thực tế và cho phép
so sánh công bằng với baseline chọn ngẫu nhiên.

\subsection{Toàn kho và thiên lệch lấy mẫu}
\label{sec:fullcatalog}

Nhiều nghiên cứu đánh giá theo cách \textit{lấy mẫu} (sampled metric): trộn sản phẩm
mục tiêu với một số ít mẫu âm (ví dụ $99$ hoặc $999$ sản phẩm) rồi xếp hạng trong
nhóm nhỏ đó. Cách này nhanh nhưng gây \textbf{thiên lệch lạc quan}: kho ứng viên càng
nhỏ, chỉ số càng cao một cách giả tạo. Hình~\ref{fig:evalproto} minh họa hiện tượng
này -- cùng một mô hình, Recall@10 giảm mạnh từ khoảng $0{,}42$ (chỉ $100$ ứng viên)
xuống $0{,}0142$ khi xếp hạng trên toàn bộ $17.000$ sản phẩm. Do đó, đồ án chủ trương
đánh giá \textbf{toàn kho (full-catalog)} để có con số trung thực, dù khắt khe hơn.

\begin{figure}[H]
\centering
\includegraphics[width=0.72\textwidth]{eval_protocol.png}
\caption{Recall@10 giảm khi số ứng viên đánh giá tăng -- minh chứng thiên lệch lấy mẫu}
\label{fig:evalproto}
\end{figure}
